% Einheit 2 – Ausklammern (terme.md Einheit 4) – Nr. 17–20
\einheitenkopf[e2][2 Ausklammern]{Einheit 2 von 5 · Ausklammern}
\zweigzeile{Hier lernst du, einen gemeinsamen Faktor auszuklammern · neu oder schon bekannt – je nach Buch · keine P10-Aufgabe · baut auf: Klammern auflösen (Einheit 1)}
\setcounter{aufgabe}{16}

\begin{aufgabe}{Ich kann einen gemeinsamen Faktor ausklammern.}
\anweisung{Klammere so viel wie möglich aus.}
\rechenplatz{2}
\begin{gleichungsraster}[1]
\gl{3x + 12} & \gl{5y + 35} \\
\gl{2x - 10} & \gl{7y + 21} \\
\gl{x^2 + 9x} & \gl{6x^2 + 15x} \\
\gl{6x + 6} & \gl{4x^2 + 8x + 12} \\
\anweisung{Klammere aus und mache die Probe durch Ausmultiplizieren.}
\gl[2]{10x^2 - 4x} & \gl[2]{12x^3 + 8x^2 - 4x} \\
\end{gleichungsraster}
\end{aufgabe}

\begin{aufgabe}{Ich finde den Fehler beim Ausklammern.}
\anweisung{Finde den Fehler, benenne ihn und korrigiere die Rechnung.}
\begin{teile}
\teil Jonas klammert aus: \rechnung{6x + 15 &= 3 \cdot (2x + 15)}
\teil Emma klammert aus: \rechnung{7x^2 + 7x &= 7x \cdot x}
\end{teile}
\end{aufgabe}

\begin{aufgabe}{Ich kann entscheiden, ob alle Glieder einen gemeinsamen Faktor haben.}
\anweisung{Haben alle Glieder einen gemeinsamen Faktor außer 1? Kreuze an, ohne zu rechnen.}
\begin{teilezwei}
\tz $4x + 9$ \quad \janein & \tz $6x + 10$ \quad \janein \\
\tz $x^2 + 3$ \quad \janein & \tz $5x^2 + 2x$ \quad \janein \\
\tz $3x + 7y$ \quad \janein & \tz $14y - 21$ \quad \janein \\
\end{teilezwei}
\end{aufgabe}

\begin{aufgabe}{Ich kann eine Sachaufgabe durch Ausklammern lösen.}
Ein rechteckiger Schulhof hat den Flächeninhalt $(8x^2 + 20x)$\,m$^2$. Er ist $4x$\,m lang.

\rechteck[seiten={4x,?,,},punkte=]{4.5}{2.5}

\begin{teile}
\teil Wie breit ist der Schulhof? Klammere dazu $4x$ aus.
\teil Für den Schulhof gilt $x = 10$. Wie lang und wie breit ist er? Prüfe mit dem Flächeninhalt.
\end{teile}
\end{aufgabe}
